第三章的身份规则已经把恢复与新购买分开。现在让首次请求 R1 和同参恢复请求 R2 同时进入服务：两者携带已经保存的完整身份
\[
 S=(1,\texttt{purchase.create.v1},E1,K1)
\]
以及 \code{paper-001}、数量 1、Q7、CNY 的原输入。

各组并发历史会就地给出初始余额，购买、订单和流水均从空记录开始。\textbf{请求在 E1 有效时进入，处理期间记录保留。}数据库沿用 PostgreSQL 18 的 Read Committed，报价 Q7 为 3000 分；后文使用这些条件检查等待、快照与提交的关系。

第三章的串行流程会先查 S，找到原结果便返回；查不到便创建购买，把身份、余额变化、订单、流水和原结果共同提交。事务把这一组修改安排为共同成立或共同回滚。现在让 R1 和 R2 同时进入服务：它们可能同时查到无记录，各自开始一笔完整购买。需要补上的条件，是两次请求都经过同一个有约束的受理步骤，且只有取得该身份的事务才继续扣款。
